% Guía UNSAAC EPIIS — transcripción del contenido de las capturas + anexo fotográfico
\documentclass[a4paper,11pt]{book}
\usepackage[utf8]{inputenc}
\usepackage[T1]{fontenc}
\usepackage[spanish]{babel}
\usepackage[margin=22mm]{geometry}
\usepackage{graphicx}
\usepackage{fancyhdr}
\usepackage{hyperref}
\usepackage{enumitem}
\usepackage{booktabs}

\hypersetup{hidelinks, pdftitle={Guía elaboración Plan de Tesis UNSAAC (transcripción)}}

\pagestyle{fancy}
\fancyhf{}
\fancyfoot[C]{\thepage}
\renewcommand{\headrulewidth}{0pt}

\title{\textbf{Guía para elaborar el Plan de Tesis}\\[0.3em]
\large Escuela Profesional de Ingeniería Informática y de Sistemas --- UNSAAC\\[0.5em]
\normalsize Transcripción del contenido a partir de las capturas \texttt{20160921\_*.jpg}}
\author{Documento compilado para uso académico}
\date{\today}

\begin{document}
\frontmatter
\maketitle
\tableofcontents
\mainmatter

\chapter*{Nota sobre esta versión}
\addcontentsline{toc}{chapter}{Nota sobre esta versión}
El texto de los capítulos siguientes reconstruye, en español, el contenido legible en las fotografías de la guía oficial. Las listas extensas (verbos) se han tomado de las descripciones fieles a las capturas; para citas normativas definitivas conviene contrastar con el documento original de la escuela. Al final del PDF se incluye el \textbf{anexo fotográfico} con las imágenes en orden.

% --- Contenido por secciones alineado al orden de las capturas ---

\chapter{Título}
\label{cap:titulo}

Un buen título debe ser breve, preciso y conciso. Debe presentar al lector o revisor las variables centrales del estudio; esas variables constituyen las palabras clave para clasificación e indexación.

\textbf{Ejemplo:} \textit{Aportes al reconocimiento de locutores para su integración en la inteligencia ambiental.}

\subsection*{Pautas para el título del proyecto}
\begin{enumerate}[leftmargin=*]
  \item Construido con palabras clave de la revisión de literatura.
  \item Debe mostrar la respuesta al problema planteado.
  \item Debe responder \textbf{¿Qué?} y \textbf{¿Para qué?}
  \item Refleja la novedad o innovación del proyecto o investigación.
  \item Alineado con el objetivo general y, por tanto, con las conclusiones.
  \item Encuadra los conceptos necesarios del marco teórico.
  \item Delimita el marco teórico del proyecto.
  \item Inicia con un verbo (acción) que refleje el alcance de la investigación.
\end{enumerate}

\subsection*{Verbos en investigación}
No todos los verbos sirven para título u objetivos de pregrado. Hay \textbf{verbos prohibidos} (subjetivos, poco medibles o de actividad sin referencia científica clara) y \textbf{verbos correctos} (medibles, verificables).

\textbf{Verbos prohibidos (extracto de la guía):} Abordar, Abrirse, Adaptar, Analizar, Armonizar, Asimilar, Atender, Buscar, Calcular, Citar, Clasificar, Colaborar, Comparar, Concientizar, Confrontar, Conocer, Considerar, Consultar, Contestar, Contrastar, Dar a conocer, Definir, Descubrir, Detectar, Difundir, Dirigir, Discriminar, Disfrutar, Distinguir, Ejercitar, Empatizar, Enlistar, Entender, Entrevistar, Enunciar, Escribir, Escuchar, Examinar, Experimentar, Exponer, Identificar, Imaginar, Informar, Iniciar, Interactuar, Interrelacionar, Investigar, Jerarquizar, Leer, Mostrar, Nombrar, Observar, Organizar, Parafrasear, Participar, Percibir, Presentar, Promover, Proponer, Realizar, Reconocer, Redactar, Reflexionar, Relacionar, Responder, Responsabilizar, Resumir, Revisar, Sensibilizar, Visualizar.

\textbf{Verbos correctos (extracto de la guía):} Aprobar, Adoptar, Aportar, Arriesgar, Asesorar, Asumir, Caracterizar, Comparar, Comprometerse, Construir, Coordinar, Crear, Criticar, Cuestionar, Decidir, Delegar, Desarrollar, Determinar, Diagnosticar, Dimensionar, Diseñar, Ejercer, Elaborar, Elegir, Emprender, Estimar, Estructurar, Evaluar, Formular, Generar, Guiar, Identificar, Innovar, Juzgar, Optar, Posicionarse, Proyectar, Simular, Sintetizar, Solucionar, Supervisar, Testimoniar, Transferir, Transformar, Transmitir, Trascender, Valorar, Valuar.

\textit{Pie de página en original: Guía para elaborar el Plan de Tesis --- EPIIS --- UNSAAC.}

\chapter{Resumen}
\label{cap:resumen}

\subsection*{Indicaciones}
\begin{itemize}[leftmargin=*]
  \item Debe incluir: tema de investigación, importancia, objetivos planteados, estrategias metodológicas y posibles aplicaciones.
  \item Redactarse en \textbf{un solo párrafo}, sin sangría ni saltos de línea.
  \item Extensión máxima \textbf{400 palabras}; Arial 12, interlineado sencillo (según guía impresa).
\end{itemize}

\subsection*{Ejemplo diagramado (tema calentamiento global)}
Párrafo de ejemplo sobre modelo climático y satélite ABC, con etiquetas: tema y antecedentes; problema; propuesta; objetivo; estrategia; importancia. \textit{Fuente en guía: adaptado de Kallestinova (2011).}

\subsection*{Ejemplo completo (informática)}
Párrafo sobre necesidad de análisis de datos, cómputo paralelo, paralelización, visión artificial, OCR con KNN y SVM, smartphones y SBC, y medición por tiempos de ejecución y consumo eléctrico. \textit{Fuente en guía: adaptado de Sandín Carral (2015).}

\chapter{Antecedentes y problema de investigación}
\label{cap:antecedentes}

\section{Antecedentes o esbozo del estado del arte}
\begin{itemize}[leftmargin=*]
  \item Extensión sugerida: 2 a 4 páginas.
  \item Debe reflejar el estado actual del conocimiento y avances en el área; sirve como modelo para investigación futura y para comparar el trabajo con estudios previos con variables u objetivos similares.
  \item Incorporar aportes teóricos de autores y especialistas.
  \item La discusión bibliográfica debe centrarse en \textbf{cómo} distintos enfoques, métodos y modelos han abordado el problema.
  \item La revisión debe ser \textbf{critica y comparativa}, mostrando ventajas y desventajas de soluciones actuales.
  \item Objetivo: describir el nivel de desarrollo del tema y el papel del contexto.
\end{itemize}

\textbf{Preguntas guía:} ¿Condiciones del contexto? ¿Cómo se ha tratado el tema? ¿Tendencias? ¿Avances? ¿Nuevas contribuciones o líneas? ¿Qué aportes sirven al presente trabajo?

\section{Problema de investigación}
\begin{itemize}[leftmargin=*]
  \item Se espera descripción amplia y detallada (mínimo sugerido: tres páginas).
  \item La problematización surge ante una necesidad concreta, vacío de conocimiento o contradicción entre enfoques.
  \item Describe una realidad donde se relacionan variables o factores (causas y consecuencias).
  \item Debe cerrar con una \textbf{pregunta de investigación} como síntesis.
  \item Propósito: describir hechos y acontecimientos y delimitar lo que se estudiará.
  \item En esta guía el planteamiento debe formularse como \textbf{pregunta}.
\end{itemize}

\chapter{Planteamiento del problema, justificación (continuación)}
\label{cap:planteamiento}

\section{Planteamiento del problema}
Incluir diagnóstico situacional, relevancia para la sociedad y situación actual (internacional, nacional, regional). Indicar qué se investiga, su naturaleza y características. Puede iniciar con interrogantes (¿En qué medida?, ¿Dónde?, ¿Cuándo?, ¿Cómo?, ¿Por qué?, ¿Qué efecto?). Delimitar tiempo, espacio y población (autores citados en guía: Hernández-Sampieri, Arias, Bernal, Intrigo).

\textbf{Ejemplo general:} ¿Qué genes inactivados en la bacteria \textit{Helicobacter pylori} confieren resistencia al antibiótico claritromicina?

\subsection*{Ejemplos en sistemas informáticos (tabla resumida)}
\begin{itemize}[leftmargin=*]
  \item \textbf{Oración:} Se desconocen los procesos administrativos y contables de la empresa X.\\
        \textbf{Pregunta:} ¿Cuáles son los procesos administrativos y contables en uso en la empresa X?\\
        \textbf{Preguntas específicas:} estructura de la empresa; quién y cuándo accede al sistema; estructura del proceso contable; demandas o usos del sistema.
  \item \textbf{Oración:} Falta información sobre sectores mejorables del sistema informático de la empresa Y.\\
        \textbf{Pregunta:} ¿Qué factores son mejorables en el sistema informático de la empresa Y?
  \item \textbf{Oración:} Requerimiento de un nuevo sistema que reemplace al obsoleto en la empresa Z.\\
        \textbf{Pregunta:} ¿Cómo debería ser un nuevo sistema informático para Z que reemplace al antiguo?
\end{itemize}

\textit{Objetivo de la sección: claridad sobre cuál es el problema, su importancia y qué se hará en general.}

\section{Justificación}
Fundamenta la importancia del problema y la necesidad del trabajo práctico para aportar solución. Expone razones (impacto, beneficiarios, aporte teórico, cambios, relevancia social).

\textbf{Preguntas guía:}
\begin{itemize}[leftmargin=*]
  \item \textbf{Conveniencia:} ¿Qué tan conveniente es la investigación? ¿Para qué sirve?
  \item \textbf{Relevancia social:} ¿Impacto social? ¿Quiénes se benefician?
  \item \textbf{Implicaciones prácticas:} ¿Ayuda a resolver un problema real? ¿Trascendencia para otros problemas prácticos?
  \item \textbf{Valor teórico:} ¿Llena vacíos, apoya teoría u orienta nuevas hipótesis?
  \item \textbf{Utilidad metodológica:} ¿Nuevos instrumentos, conceptos o métodos experimentales?
  \item \textbf{Consecuencias del estudio:} repercusiones positivas o negativas en la comunidad.
\end{itemize}

\chapter{Objetivos y verbos}
\label{cap:objetivos}

\section{Objetivos}
Un objetivo es una afirmación clara de lo que se persigue; responde \textbf{¿Qué?}, \textbf{¿Para qué?} y \textbf{¿A través de qué?}

Características: concretos y verificables; no son metas lejanas ni meras consecuencias; la investigación se evalúa por su logro.

Evitar objetivos triviales o confundir objetivos con métodos (ej.\ ``estudiar satisfacción mediante entrevistas'' como objetivo único).

Requisitos: \textbf{factible}, \textbf{verificable}, \textbf{preciso}, orientado al \textbf{problema}.

\subsection*{Secuencia sintagmática (tabla guía)}
Columnas: Verbo --- Fenómeno --- Subfenómeno (opcional) --- Para\ldots\ (finalidad). Verbos ejemplo: Establecer, Averiguar, Identificar, Recopilar, Investigar, Revelar, Descubrir, Indagar, Inquirir, Pesquisar, Registrar, Buscar\ldots\ Adaptado de Romelia Rodríguez (2008).

\section{Verbos para objetivos generales y específicos}
Listas numeradas en la guía (extracto): \textit{Objetivo general:} Analizar, Calcular, Categorizar, Comparar, Crear, Definir, Desarrollar, Diseñar, Evaluar, Formular, Generar, Identificar, Proponer, Reconstruir, Valuar\ldots

\textit{Objetivos específicos:} Advertir, Basar, Conceptualizar, Contrastar, Determinar, Describir, Enumerar, Especificar, Estimar, Interpretar, Operacionalizar, Organizar, Registrar, Relacionar, Sintetizar, Sugerir\ldots

\section{Taxonomía de Bloom (niveles cognitivos)}
\begin{itemize}[leftmargin=*]
  \item \textbf{Conocimiento:} definir, describir, enumerar, identificar, memorizar\ldots
  \item \textbf{Comprensión:} explicar, interpretar, resumir, traducir\ldots
  \item \textbf{Aplicación:} aplicar, calcular, ejecutar, implementar\ldots
  \item \textbf{Análisis:} analizar, comparar, diferenciar, examinar\ldots
  \item \textbf{Síntesis:} crear, diseñar, elaborar, formular, proponer\ldots
  \item \textbf{Evaluación:} evaluar, juzgar, criticar, validar\ldots\ (la guía marca como \textit{no recomendable} ciertos verbos vagos: apreciar, comprender, desear, disfrutar, entender, dominar, familiarizarse, interesarse en\ldots).
\end{itemize}

\section{Ejemplos de objetivos (guía)}
\textbf{Ejemplo 1} (Richards, 2011): Objetivo general sobre genes inactivados en \textit{H.\ pylori}; específicos sobre screening y cepas naturales.

\textbf{Ejemplo 2} (Sandín Carral, 2015): Paralelismo y OCR; específicos sobre algoritmos, Android, plataformas, análisis comparativo.

\textbf{Ejemplo 3:} Logística inversa PET; específicos sobre factores de reciclaje, impacto, finalidad de reducción de desechos, parámetros para propuesta.

\section{Tipos de investigación y verbos asociados}
Tabla guía (resumen):

\begin{center}
\small
\begin{tabular}{@{}p{2.2cm}p{3.5cm}p{3cm}p{3.8cm}@{}}
\toprule
\textbf{Tipo} & \textbf{Propósito} & \textbf{Obj.\ gral.} & \textbf{Obj.\ específicos} \\
\midrule
Descriptiva & Características, cómo es & Describir, Diagnosticar & Identificar, estudiar, detectar, categorizar\ldots \\
Comparativa & Manifestación en grupos & Comparar & Contrastar, diferenciar\ldots \\
Analítica & Elementos del fenómeno & Analizar & Criticar, descomponer\ldots \\
Explicativa & Causas, por qué & Explicar & Deducir, inferir, relacionar\ldots \\
Proyectiva & Diseño o propuesta & Proponer & Formular, diseñar, planear\ldots \\
Experimental & Relación entre variables & Confirmar & Verificar, comprobar, demostrar\ldots \\
Evaluativa & Logro de objetivos & Evaluar & Valorar, estimar, juzgar\ldots \\
\bottomrule
\end{tabular}
\end{center}

\chapter{Hipótesis}
\label{cap:hipotesis}

La hipótesis es \textbf{opcional} según la naturaleza del trabajo. Es explicación tentativa del fenómeno, formulada como proposición; respuesta provisional a preguntas de investigación (Hernández, Fernández y Baptista, 2010).

Características:
\begin{enumerate}[leftmargin=*]
  \item Referirse a situación real y ser contrastable en un universo y contexto definidos.
  \item Variables comprensibles, precisas y concretas.
  \item Relación entre variables clara y plausible.
  \item Términos observables y medibles.
  \item Deben relacionarse con técnicas disponibles para probarlas.
\end{enumerate}

\subsection*{Clasificación (notación)}
Hipótesis de investigación (Hi); nula (Ho); alternativa (Ha); estadística (traducción formal de Hi).

Tipos: descriptivas de valor pronosticado; correlacionales; diferencia de grupos; causales.

\chapter{Alcances, limitaciones y marco teórico (esbozo)}
\label{cap:alcance}

\section{Alcances y limitaciones}
Presentar en forma de puntos, redacción clara y detallada.

\textbf{Alcances:} dependen del tema y del alcance deseado. En estudios cuantitativos: exploratorio, descriptivo, correlacional, explicativo.

\textbf{Limitaciones:} actividades; recursos; tiempo; espacio o territorio.

\section{Esbozo del marco teórico}
Presentar brevemente temas que se desarrollarán, derivados de palabras clave y del título (Sautu et al., 2005). El marco teórico es el conjunto de conceptos de distintos niveles de abstracción que orientan la forma de aprehender la realidad; incluye paradigma, teoría general (proposiciones interrelacionadas que explican procesos y fenómenos), etc. En esta etapa basta listar temas, no desarrollarlos completamente.

\chapter{Metodología, resultados esperados y contribuciones}
\label{cap:metodologia}

\section{Metodología}
Relación amplia y detallada; extensión según el método elegido. La metodología es el conjunto de procedimientos racionales y ordenados para alcanzar objetivos. Redactar el procedimiento en párrafos.

\textbf{Subsección ``Sistema/modelo/método desarrollado'':} describir con detalle técnico cómo se abordó el problema: métodos, modelos, teorías, técnicas, decisiones, supuestos y diferencias frente a otros enfoques. Distinguir lo técnicamente desarrollado de su aplicación en un contexto.

Ejemplo tabular (reingeniería de software / empresa marmolera): fases de análisis de metodologías OAR y modelo cíclico, combinación para la industria, inventario de componentes, restauración de documentos. \textit{Fuente en guía: Urquizo B.\,E.\ et al.\ (2015).}

\textit{Objetivo:} responder \textbf{qué} se hizo, \textbf{cómo} se hizo y en qué puede diferir de alternativas.

\section{Resultados esperados (a priori)}
Lo que se espera obtener al finalizar el proyecto o investigación.

\section{Contribuciones originales esperadas}
Redactar contribuciones originales previstas al concluir, indicando en qué aspecto y áreas del conocimiento se enriquecen.

\chapter{Impacto social, índice, cronograma, presupuesto y referencias}
\label{cap:impacto}

\section{Impacto social esperado}
Efectos sobre la población; identificar \textbf{región geográfica} o \textbf{sector poblacional} beneficiado.

\section{Índice tentativo}
Índice detallado factible para desarrollar la investigación.

\section{Cronograma}
Presentar cronograma; se menciona \textbf{diagrama de Gantt} (Bernal, 2010) como base ordenada para estudiante y asesor.

\section{Presupuesto}
Recursos necesarios; estimación de inversión (materiales, viajes, etc.); indicar fuentes de financiamiento. Tabla guía: Componentes; Costos (S/.); Fuente de financiamiento. Rubros: materiales, equipos, software, viáticos, transporte, gastos administrativos, otros no previstos.

\section{Referencias y citas bibliográficas}
Cumplir formato \textbf{APA} (en la guía impresa se menciona la guía vigente en la fecha del documento). Todo material usado debe citarse. Ejemplos en guía: libros (Comyns, 2007); artículos (Henikoff); publicaciones periódicas (Nawrot et al.); DOI; tesis; informes técnicos; documentos electrónicos.

\chapter{Formato de presentación y normas editoriales}
\label{cap:formato}

\section{Portada (Anexo A)}
Logo UNSAAC; nombre completo de la universidad; título del plan; nombre y código del estudiante; asesor(a) y co-asesor(a); Perú, mes y año; Escuela Profesional de Ingeniería Informática y de Sistemas; Facultad de Ingeniería Eléctrica, Electrónica, Informática y Mecánica; nombre completo de la universidad (pie).

\section{Márgenes y papel}
Tamaño \textbf{A4}. Márgenes superior, inferior y derecho \textbf{2{,}54 cm}; izquierdo \textbf{3{,}81 cm}.

\section{Tipografía}
Arial o Times New Roman \textbf{12 pt}; interlineado \textbf{1{,}5}; texto \textbf{justificado}.

\section{Citas en el texto (APA)}
Sistema autor-fecha; citar al usar ideas o textos ajenos; incluir apellido, año y página en citas textuales directas.

\section{Numeración}
Todas las páginas numeradas abajo a la derecha; excepción portada e índices. Índices en romanos. La portada cuenta como i pero sin imprimir número.

\section{Título del plan}
Puede ir en mayúsculas 12 pt; acentos correctos; nombres científicos en cursiva si aplica.

\section{Subdivisiones}
Mayúsculas 12 pt; negrita posible; alineadas a la izquierda.

\section{Figuras}
Autoevidentes; 3--4 grupos de datos por figura; escalas y símbolos claros; \textbf{título} arriba a la izquierda en 10 pt cursiva; numeración consecutiva e independiente de tablas; indicar fuente si no es propia. Ejemplo en guía: Figura 1 sobre algoritmo genético simple (Goldberg, 1989).

\section{Tablas y cuadros}
Título superior en \textbf{10 pt} (puede ir en negrita); fuente precisa; numeración independiente de figuras; citar fuente si no es propia. Ejemplo: Tabla 1 error PCA locutores (Alcocer et al., 2004).

\section{Anexos}
Lista A: elementos de portada (misma lista que en sección de portada).

% --- Anexo fotográfico ---
\backmatter
\chapter*{Anexo: capturas originales}
\addcontentsline{toc}{chapter}{Anexo: capturas originales}
Las siguientes páginas reproducen las fotografías \texttt{20160921\_150046.jpg} a \texttt{20160921\_150257.jpg} en orden, para cotejo con esta transcripción.

\newcommand{\fotoguia}[1]{%
  \thispagestyle{fancy}%
  \begin{center}
    \includegraphics[width=\linewidth,height=\textheight,keepaspectratio]{#1}\\[0.5em]
    \texttt{\detokenize{#1}}
  \end{center}
  \clearpage
}

\fotoguia{20160921_150046.jpg}
\fotoguia{20160921_150057.jpg}
\fotoguia{20160921_150111.jpg}
\fotoguia{20160921_150121.jpg}
\fotoguia{20160921_150129.jpg}
\fotoguia{20160921_150137.jpg}
\fotoguia{20160921_150147.jpg}
\fotoguia{20160921_150155.jpg}
\fotoguia{20160921_150204.jpg}
\fotoguia{20160921_150213.jpg}
\fotoguia{20160921_150224.jpg}
\fotoguia{20160921_150231.jpg}
\fotoguia{20160921_150240.jpg}
\fotoguia{20160921_150246.jpg}
\fotoguia{20160921_150257.jpg}

\end{document}
